\section{Background and Related Work}\label{sec:related}

\subsection{Scheduling in Kubernetes and in multi-cluster systems}\label{sec:related-k8s}

Kubernetes inherited its scheduling model from cluster managers such as Borg, which places tasks on machines using admission control and resource packing \cite{verma2015borg}. The Kubernetes scheduler selects a node for each pod through a framework of plugins at defined extension points (filtering, scoring, reserving, binding and others), and the scores of the scoring plugins are combined using configurable weights \cite{k8s_scheduling_framework}. Surveys classify the large body of work that extends this scheduler by objective, workload type and environment \cite{carrion2022kubernetes,rejiba2022custom,senjab2023survey}. The most recent review covers 2023 to early 2026 and includes scheduling scopes from cloud to edge \cite{alhakimi2026enhancing}.

Placement across several independent clusters is a separate problem, because each cluster keeps its own control plane and its own node-level scheduler. KubeFed, the earlier federation project of the Kubernetes community, was archived in April 2023 \cite{kubefed_repo}. Liqo shares resources between clusters by representing a remote cluster as a virtual node of the local one \cite{iorio2023liquid}. Oakestra, an orchestrator for edge infrastructure that does not build on Kubernetes, uses a hierarchy in which a root component first selects a cluster and the selected cluster then places the task \cite{bartolomeo2023oakestra}. Karmada follows the same two-level decomposition while remaining Kubernetes-native: a central control plane selects member clusters, and the kube-scheduler of each member cluster selects nodes \cite{karmada_docs_architecture}.

\subsection{Karmada}\label{sec:related-karmada}

A workload submitted to Karmada is an ordinary Kubernetes resource template, for example a Deployment. A PropagationPolicy states the placement rules, such as cluster affinity, spread constraints and the replica scheduling type. The control plane combines template and policy into a ResourceBinding, the Karmada scheduler writes the selected clusters into that binding, and controllers then create the corresponding objects in the member clusters \cite{karmada_docs_architecture}. The scheduler can be extended with filter and score plugins compiled into a custom scheduler binary, and several scheduler instances can run side by side, each serving the workloads whose policy names it \cite{karmada_docs_custom_scheduler}. Rebalancing of workloads that are already placed is a separate, explicitly triggered mechanism \cite{karmada_docs_rebalancer}. An empirical study of Karmada has examined how applications are reconfigured after cluster failures or resource shortage \cite{patino2025reconfiguration}. The parts of the scheduler that matter for this paper are analysed in \cref{sec:karmada-analysis}.

\subsection{Network-aware scheduling within one cluster}\label{sec:related-node}

Santos et al.\ extended the Kubernetes scheduler to account for the round-trip time between nodes for fog applications \cite{santos2019networkaware}. Their later framework, Diktyo, places dependent microservices using an application dependency graph with latency and bandwidth requirements \cite{santos2023diktyo}; its plugins are distributed in the Kubernetes scheduler-plugins project \cite{schedulerplugins_networkaware}. NetMARKS scores nodes using metrics collected by the Istio service mesh \cite{wojciechowski2021netmarks}, and NAS uses latency and bandwidth observed through a Cilium eBPF service mesh \cite{ragavan2025nas}. Centofanti et al.\ argue for measuring latency at the application layer, as perceived by end users, instead of between cluster nodes \cite{centofanti2023latency}. Marchese and Tomarchio combine a network-aware scheduler plugin with a descheduler that moves pods when the network state or the traffic between microservices changes \cite{marchese2022ccgrid}, and they treat resource load in a separate strategy \cite{marchese2025loadaware}. Rossi et al.\ combine elasticity with network-aware placement for nodes that are geographically distributed but managed by one control plane \cite{rossi2020gekube}.

These works establish that network state is a useful scheduling signal, but they operate on nodes within one control plane. At cluster level the available resource information is aggregated over many nodes, it is refreshed periodically rather than observed per pod, and the relevant network path runs from clients to the entry point of a cluster rather than between nodes.

\subsection{Cluster selection with resource and network information}\label{sec:related-cluster}

Several works already combine resource and latency information when selecting clusters. Tabatabaei et al.\ build a service orchestration architecture on Karmada in which a filter plugin hands the candidate clusters to an external integer linear programming solver; the objective weighs peak CPU load against latency, and the authors state that the scoring mechanisms of Karmada are not yet included in their flow \cite{tabatabaei2025cnf}. Their simulation models latency with fixed round-trip values without running service traffic, and the authors name the focus on initial placement as a limitation. CosmoSpan selects a cluster and a node jointly through a multi-criteria scheduler that binds through the Karmada API, together with deschedulers that react to changes \cite{bregar2026cosmospan}. Its placement score covers CPU, memory, node uptime, bandwidth, latency and fault tolerance; it is evaluated mainly in simulation, and the Karmada scheduler is not part of the comparison because no simulation framework exists for it \cite{bregar2026cosmospan}. LAIS is a multi-cluster scheduler driven by end-to-end latency measurements that either meets a latency constraint or chooses the cluster with the lowest latency, and it is compared with the default Kubernetes scheduler \cite{chiaro2024lais}. Earlier, mck8s provided controllers on top of KubeFed that place, scale and burst applications across geo-distributed clusters and was evaluated on the Grid'5000 testbed \cite{tamiru2021mck8s}. The CODECO preprint ranks clusters using CPU, memory, bandwidth, latency and service class over Open Cluster Management \cite{sofia2026codeco}. Formal models for placement across federated domains also exist \cite{faticanti2023locality,khan2025latency}, as do learned policies based on deep reinforcement learning for multi-cluster scheduling \cite{huang2026kmps,jiang2025drkc}.

\Cref{tab:related} summarises how the closest works relate to the problem of this paper.

\begin{table}[t]
\caption{Cluster-level placement approaches that use resource and network information}\label{tab:related}
\centering\small\setstretch{1.0}
\begin{tabularx}{\textwidth}{@{}p{2.3cm}LLL@{}}
\toprule
\textbf{Work} & \textbf{Integration with orchestrator} & \textbf{Signals and decision} & \textbf{Missing-data handling, evaluation} \\
\midrule
Tabatabaei et al.\ \cite{tabatabaei2025cnf} & Karmada filter plugin calling an external ILP solver & Weighted peak CPU load and latency & Not specified; simulation with fixed RTT values, testbed deployment time \\
CosmoSpan \cite{bregar2026cosmospan} & Own scheduler; binds through the Karmada API & Multi-criteria score; latency as a filter & Not specified; KWOK simulation, multi-region overhead study \\
LAIS \cite{chiaro2024lais} & Multi-cluster scheduler with latency intents & End-to-end latency (hard or soft intent) & Not specified; comparison with the default Kubernetes scheduler \\
mck8s \cite{tamiru2021mck8s} & Controllers on KubeFed (archived) & Placement, scaling and bursting policies & Not specified; Grid'5000 testbed \\
CODECO \cite{sofia2026codeco} & Open Cluster Management & CPU, memory, bandwidth, latency, service class & Not specified; preprint without quantitative placement evaluation \\
This paper & Karmada filter and score plugins; default stages kept & Resource headroom and client-to-cluster RTT; capacity as a hard constraint & Explicit validity and fallback rules; design only, validation planned \\
\bottomrule
\end{tabularx}
\end{table}

Two observations follow from this comparison. First, none of the works defines how a decision is made when network measurements are missing or out of date, although measurements in a real system are periodic and can fail. Second, the works that build on Karmada either delegate the decision to an external solver or bypass the Karmada cluster scheduler, so the interaction of a resource- and network-aware objective with Karmada's own filtering, scoring and selection stages has not been described. This paper addresses both points at the design level. The empirical comparison against the default Karmada scheduler is left to the validation proposed in \cref{sec:validation}.
